% Part: first-order-logic
% Chapter: beyond
% Section: second-order-logic

\documentclass[../../../include/open-logic-section]{subfiles}

\begin{document}

\olfileid{fol}{byd}{sol}

\olsection{Logika Orde Kapindho}

Basa logika orde kapindho ngidini kuantifikasi
ora mung tumrap !!a{domain} individu, nanging
uga tumrap relasi ing !!{domain} kasebut.
Saka basa logika tataran kapisan $\Lang{L}$, kanggo
saben $k$ ditambahake !!{variable}s $R$ kang
nglimputi relasi $k$-ary, lan kuantifikasi
tumrap variabel-variabel iku diidinake.
Yen $R$ iku !!a{variable} kanggo relasi
$k$-ary, lan $t_1$, \dots,~$t_k$ iku
term biasa (orde kapisan), mula
$\Atom{R}{t_1,\dots,t_k}$ iku !!{formula}
atomik. Saliyane iku, himpunan !!{formula}s
ditemtokake kaya ing logika tataran kapisan,
kanthi klausa tambahan kanggo kuantifikasi
orde kapindho. Gatekna yen !!{identity}
mung ana kanggo term orde kapisan: yen
$R$ lan $S$ iku !!{variable}s relasi
kanthi aritas padha~$k$, kita bisa
nemtokake $\eq[R][S]$ minangka cekakan
kanggo
\[
\lforall[x_1 \dots][\lforall[x_k][(\Atom{R}{x_1, \dots, x_k} \liff
  \Atom{S}{x_1, \dots, x_k})]].
\]

Aturan logika orde kapindho mung nggedhekake
aturan kuantor supaya nyakup variabel
orde kapindho kang anyar. Nanging ing kene
kita kudu rada ngati-ati njlentrehake
sesambungane variabel-variabel iki karo
!!{predicate}s saka~$\Lang{L}$, lan
luwih umum karo !!{formula}s saka~$\Lang{L}$.
Paling ora, variabel relasi dianggep term,
mula ana inferensi kanthi wujud
\[
!A(R) \Proves \lexists[R][!A(R)]
\]
Nanging yen $\Lang{L}$ iku basa aritmetika
kanthi simbol relasi tetep~$<$, kita uga
bakal nganggep inferensi ing ngisor iki sah:
\[
x < y \Proves \lexists[R][\Atom{R}{x,y}]
\]
utawa kanggo !!{formula}~$!A$ tartamtu,
\[
!A(x_1, \dots, x_k) \Proves \lexists[R][\Atom{R}{x_1,\dots,x_k}]
\]
Luwih umum, kita bisa ngidini inferensi
kanthi wujud
\[
\Subst{!A}{\lambd[\vec x][!B(\vec x)]}{R} \Proves
\lexists[R][!A]
\]
ing ngendi $\Subst{!A}{\lambd[\vec x][!B(\vec x)]}{R}$
nuduhake asil ngganti saben !!{formula}
atomik kang variabel relasine katon bebas, kanthi wujud
$\Obj{R}{t_1,\dots,t_k}$ ing~$!A$
nganggo $!B(t_1, \dots, t_k)$. Cathetan cakupan SOL6-F013: panggantos iki mung kanggo occurrence bebas variabel relasi; variabel kaiket kudu diwenehi jeneng anyar yen perlu supaya variabel bebas formula panggantos ora kaiket kanthi ora sengaja.
Aturan kang pungkasan iki setara karo
anane {\em skema komprehensi}, yaiku
aksioma kanthi wujud
\[
\lexists[R][\lforall[x_1, \dots, x_k][(!A(x_1, \dots, x_k) \liff
\Atom{R}{x_1, \dots, x_k})]],
\]
siji kanggo saben !!{formula}~$!A$
ing basa orde kapindho, kanthi
$R$ ora dadi variabel bebas.
(Gladhen: tuduhna yen $R$ diidinake
katon bebas ing~$!A$, skema iki ora konsisten!)

Nalika ahli logika nyebut "aksioma logika
orde kapindho", lumrahe kang dimaksud
yaiku pangembangan paling cilik saka
logika tataran kapisan kanthi aturan kuantor
orde kapindho lan skema komprehensi.
Nanging subsistem aksioma lan aturan iki
kang luwih ringkih uga asring narik
kawigaten kanggo ditliti. Upamane,
ing wujud paling umume, skema aksioma
komprehensi iku \emph{impredikatif}:
skema iki ngidini kita nyatakake anane
relasi $\Atom{R}{x_1, \dots, x_k}$
kang "ditemtokake" dening !!a{formula}
kanthi kuantor orde kapindho; dene
kuantor-kuantor iki nglimputi himpunan
kabeh relasi kaya mangkono---himpunan
kang uga ngemot $R$ dhewe!{}
Ing wiwitan abad kaping rong puluh,
salah siji tanggapan kang umum marang
paradoks Russell yaiku nganggep
definisi kaya mangkono minangka
sumber masalah, banjur nyingkiri
definisi iku nalika ngembangake
dhasar matematika. Yen kuantor
orde kapindho ora diidinake ing
!!{formula}~$!A$, kita oleh
wujud komprehensi {\em predikatif},
kang rada luwih ringkih.

Saka sudut pandang semantik,
!!{structure} orde kapindho bisa
dianggep dumadi saka !!{structure}
orde kapisan kanggo basa kasebut,
lan himpunan relasi ing !!{domain}
kang dadi cakupan kuantor orde
kapindho (luwih tepat, kanggo
saben $k$ ana himpunan relasi
kanthi aritas~$k$). Mesthi,
yen komprehensi dilebokake ing
sistem !!{derivation}, banjur
diperlokake syarat tambahan:
ana cukup relasi ing "perangan
orde kapindho" kanggo nyukupi
aksioma komprehensi---yen ora,
sistem !!{derivation} iku ora sah!{}
Salah siji cara gampang kanggo
njamin anane cukup relasi yaiku
njupuk perangan orde kapindho
kang ngemot \emph{kabeh} relasi
ing perangan orde kapisan.
!!a{structure} kaya mangkono
diarani \emph{pepak}, lan ing
sawijining pangertèn pancen
iku "!!{structure} kang dikarepake"
kanggo basa kasebut. Yen kita
mbatesi perhatian mung marang
!!{structure}s pepak, kita
oleh kang diarani semantik
orde kapindho \emph{pepak}.
Ing kahanan iki, nemtokake
!!{structure} mung mbutuhake
nemtokake perangan orde kapisan,
amarga isi perangan orde kapindho
wis ditemtokake kanthi implisit.

Dadi, nalika ngrembug logika
orde kapindho ana sawetara
teges kang bisa beda. Kanggo
sistem !!{derivation}, kang
dimaksud bisa salah siji saka
\begin{enumerate}
\item Sistem !!{derivation}
  orde kapindho "minimal",
  bebarengan karo sawetara
  aksioma komprehensi.
\item Sistem !!{derivation}
  orde kapindho "standar",
  kanthi komprehensi pepak.
\end{enumerate}
Kanggo semantike, kang
digatekake bisa salah siji saka
\begin{enumerate}
\item Semantik "ringkih",
  ing ngendi !!a{structure}
  dumadi saka perangan orde
  kapisan lan perangan orde
  kapindho kang cukup gedhe
  kanggo nyukupi aksioma
  komprehensi.
\item Semantik orde kapindho
  "standar", kang mung
  nggatekake !!{structure}s pepak.
\end{enumerate}
Yen ahli logika ora nyebutake
sistem !!{derivation} utawa
semantik kang dimaksud,
lumrahe nuduhake butir kapindho
ing saben dhaptar. Kauntungan
nggunakake semantik iki yaiku,
kaya kang bakal kita deleng,
kita bisa njlentrehake akeh
struktur matematika lumrah
kanthi kategoris; bebarengan
iku, sistem !!{derivation}
cukup kuwat lan sah kanggo
semantik iki. Kakurangane,
sistem !!{derivation}
\emph{ora} jangkep kanggo
semantik kasebut; nyatane,
\emph{ora ana} sistem
!!{derivation} kang bisa
ditemtokake kanthi efektif
kang jangkep kanggo semantik
orde kapindho pepak. Dene
kanggo semantik kang luwih
ringkih, kaya kang bakal
kita deleng, sistem
!!{derivation} iku
\emph{jangkep}; iki ateges
yen sawijining ukara ora
bisa dibuktekake, banjur
ana \emph{sawijining}
!!{structure}, ora mesthi
sing pepak, kang ndadekake
ukara iku ora bener.

Basa logika orde kapindho
sugih banget. Relasi unari
bisa kita samakake karo
subhimpunan saka !!{domain},
mula kita uga bisa
nguantifikasi himpunan-himpunan
kasebut; upamane, induksi
kanggo wilangan asli bisa
dinyatakake kanthi siji aksioma
\[
\lforall[R][((\Atom{R}{\Obj{0}} \land \lforall[x][(\Atom{R}{x} \lif
    \Atom{R}{x'})]) \lif \lforall[x][\Atom{R}{x}])].
\]
Yen basa aritmetika nduweni
simbol $\Obj 0, \Obj \prime, +,
\times$ lan $<$, kita bisa
nambah aksioma ing ngisor iki
kanggo njlentrehake tumindake:
\begin{enumerate}
\item $\lforall[x][\lnot x' = \Obj 0]$
\item $\lforall[x][\lforall[y][(x' = y' \lif x = y)]]$
  Cathetan koreksi sumber OLPL-084: simbol s ora kalebu
  ing dhaptar simbol basa; aksioma injektivitas penerus
  nggunakake notasi tandha petik kang wis ditemtokake.
\item $\lforall[x][(x + \Obj 0) = x]$
\item $\lforall[x][\lforall[y][(x + y') = (x + y)']]$
\item $\lforall[x][(x \times \Obj 0) = \Obj 0]$
\item $\lforall[x][\lforall[y][(x \times y') = ((x \times y) + x)]]$
\item $\lforall[x][\lforall[y][(x < y \liff \lexists[z][y = (x + z')])]]$
\end{enumerate}
Ora angel nuduhake yen aksioma
iki, bebarengan karo aksioma
induksi ing ndhuwur, menehi
gambaran kategoris marang
!!{structure}~$\Struct{N}$,
model standar aritmetika,
yen kita nggunakake semantik
orde kapindho pepak.
Kanggo sembarang
!!{structure}~$\Struct{M}$
kang nyukupi aksioma iki,
temtokna fungsi~$f$
saka $\Nat$ menyang
!!{domain} saka $\Struct{M}$
kanthi rekursi biasa ing $\Nat$,
saéngga $f(0) = \Assign{\Obj 0}{M}$
lan $f(x+1) = \Assign{\prime}{M}(f(x))$.
Kanthi induksi biasa
ing~$\Nat$ lan amarga aksioma
(1) lan~(2) bener ing~$\Struct M$,
kita weruh yen $f$ iku
!!{injective}. Kanggo nuduhake
yen $f$ iku !!{surjective},
upamakna~$P$ iku himpunan
unsur-unsur ing~$\Domain{M}$
kang kalebu ing jangkauan~$f$.
Amarga $\Struct M$ pepak,
$P$ kalebu ing !!{domain}
orde kapindho. Saka pambangunan~$f$,
kita ngerti yen $\Assign{\Obj
  0}{M}$ kalebu ing~$P$,
lan $P$~katutup tumrap
$\Assign{\prime}{M}$.
Amarga aksioma induksi bener
ing $\Struct M$ (mligine
kanggo~$P$), mula $P$~padha
karo sakabehing !!{domain}
orde kapisan saka~$\Struct M$.
Iki nuduhake yen $f$
iku !!a{bijection}.
Nuduhake yen $f$ iku
homomorfisme ora luwih angel,
kanthi nggunakake induksi
biasa ing $\Nat$ bola-bali.

Ing istilah teori himpunan,
fungsi mung jinis relasi
tartamtu; upamane, fungsi
unari~$f$ bisa disamake
karo relasi binari~$R$
kang nyukupi
$\lforall[x][\lexists![y][R(x,y)]]$.
Mula, kita uga bisa
nguantifikasi fungsi.
Kanthi semantik pepak,
kita bisa nemtokake kelas
!!{structure}s tanpa wates
minangka kelas
!!{structure}s $\Struct M$
kang nduweni fungsi injektif
saka !!{domain} saka $\Struct M$
menyang subhimpunan sejati
saka ranah iku dhewe:
\[
\lexists[f][(\lforall[x][\lforall[y][(\eq[f(x)][f(y)] \lif \eq[x][y])]]
  \land \lexists[y][\lforall[x][\eq/[f(x)][y]]])].
\]
Negasi ukara iki banjur
nemtokake kelas
!!{structure}s winates.

Kajaba iku, kita bisa
nemtokake kelas urutan rapi
kanthi nambah rumus ing
ngisor iki marang definisi
urutan linear:
\[
\lforall[P][(\lexists[x][\Atom{P}{x}] \lif \lexists[x][(\Atom{P}{x}
    \land \lforall[y][(y < x \lif \lnot \Atom{P}{y})])])].
\]
Rumus iki nyatakake yen
saben himpunan ora kosong
nduweni unsur paling cilik,
kanthi nganggep "himpunan"
padha karo "relasi unari".
Tuladha liyane, kita bisa
nyatakake sipat keterhubungan
graf kanthi kandha yen ora
ana pamisahan titik kang
ora sepele dadi perangan-
perangan kang ora kaubung:
\[
\lnot \lexists[A][(\lexists[x][A(x)] \land \lexists[y][\lnot A(y)]
  \land \lforall[w][\lforall[z][((\Atom{A}{w} \land \lnot \Atom{A}{z})
      \lif \lnot \Atom{R}{w,z})]])].
\]
Minangka gladhen liyane,
kita bisa nyoba nemtokake
kelas !!{structure}s winates
kang !!{domain}-e nduweni
cacah genep. Sing luwih
narik kawigaten, kita bisa
menehi gambaran kategoris
marang wilangan real minangka
lapangan urut kang jangkep
lan ngemot wilangan rasional.

Dadi, logika orde kapindho
luwih kuwat banget anggone
nyatakake tinimbang logika tataran
kapisan. Iku kabar becik;
saiki kabar alane. Kita wis
nyebutake yen ora ana sistem
!!{derivation} efektif
kang jangkep kanggo semantik
orde kapindho pepak. Kanthi
kauntungan lan kakurangane,
akeh sipat logika tataran kapisan
ora ana ing kene, kalebu
kompaktitas lan teorema
L\"owenheim--Skolem.

Kosok balene, yen kita gelem
ninggalake semantik orde
kapindho pepak lan nggunakake
semantik kang luwih ringkih,
sistem !!{derivation}
orde kapindho minimal jangkep
kanggo semantik iki. Kanthi
tembung liya, yen $\Proves$
diwaca "bisa dibuktekake
ing sistem minimal" lan
$\Entails$ diwaca
"ngakibatake sacara logis
ing semantik ringkih",
kita bisa nuduhake yen
kapan wae $\Gamma \Entails !A$,
mula $\Gamma \Proves !A$.
Yen kita arep nambah aksioma
komprehensi tartamtu ing
sistem !!{derivation},
kita kudu mbatesi semantik
marang struktur orde
kapindho kang nyukupi
aksioma kasebut: upamane,
yen $\Delta$ iku himpunan
aksioma komprehensi
(bisa uga kabeh), mula yen
$\Gamma
\cup \Delta \Entails !A$,
banjur $\Gamma \cup \Delta \Proves !A$.
Mligine, yen $!A$ ora
bisa dibuktekake nganggo
aksioma komprehensi kang
kita gatekake, banjur
ana model saka $\lnot !A$
kang tetep nyukupi
aksioma komprehensi kasebut.

Cara paling gampang kanggo
ndeleng yen teorema
kelengkapan bener kanggo
semantik kang luwih ringkih
yaiku nganggep logika orde
kapindho minangka logika
akeh jinis, kaya mangkene.
Siji jinis diinterpretasi
minangka !!{domain}
"orde kapisan" kang lumrah,
lan kanggo saben $k$
ana !!a{domain} "relasi
kanthi aritas $k$".
Kita njupuk basa kang
nduweni simbol relasi
bawaan "$\Atom{\Obj{true}_k}{R,x_1,\dots,x_k}$"
kang dimaksud kanggo
nyatakake yen $R$ bener
kanggo $x_1$, \dots,~$x_k$,
ing ngendi $R$ iku
variabel jinis "relasi
$k$-ary" lan $x_1$,
\dots,~$x_k$ iku objek
jinis orde kapisan.

Kanthi panyamaan iki,
semantik orde kapindho
ringkih sejatine padha
karo semantik lumrah
kanggo logika akeh jinis;
lan kita wis ngerteni
yen logika akeh jinis
bisa dilebokake ing
logika tataran kapisan. Dadi,
kanthi nggatekake
terjemahan bolak-balik,
pangertèn ringkih
logika orde kapindho
sejatine wujud logika tataran
kapisan kang nyamar,
ing ngendi !!{domain}
ngemot "objek" lan
"relasi" kang manut
aksioma kang cocog. Cathetan cakupan SOL6-F014: domain relasi ing semantik lemah dumadi saka relasi nyata, mula rong obyek jinis relasi kang duwe ekstensi padha kudu diidentifikasi. Ing encoding akeh jinis, iki diwenehake kanthi aksioma ekstensionalitas kanggo jinis relasi utawa kanthi njupuk quotient marang ekstensi kang padha.

\end{document}
